\begin{table*}[t]
\centering
\small
\setlength{\tabcolsep}{4pt}
\begin{tabular}{lllll}
\toprule
\textbf{Condition}
& \textbf{Control unit}
& \textbf{Goal network}
& \textbf{Scheduling}
& \textbf{Execution} \\
\midrule
\multicolumn{5}{l}{\textit{Baselines}} \\
\midrule
\condSinglePass
& Document
& \no
& \no
& One context \\
\condStaged~\cite{yang2022re3,wan2025cognitive}
& Stage
& \no
& Fixed stages
& One context \\
\condTaskPlanning~\cite{xiong2025beyond}
& Task node
& \no
& Adaptive tasks
& Delegated tasks \\
\midrule
\multicolumn{5}{l}{\textit{Process-level conditions}} \\
\midrule
\condFull
& Writing process
& \yes
& Adaptive processes
& Delegated roles \\
\condNoGoals
& Writing process
& \no
& Adaptive processes
& Delegated roles \\
\condFixedOrder
& Writing process
& \yes
& Fixed cycle
& Delegated roles \\
\midrule
\multicolumn{5}{l}{\textit{Diagnostic ablation}} \\
\midrule
\condSingleWriter
& Paragraph
& \yes
& Writer-decided
& One context \\
\bottomrule
\end{tabular}
\caption{Experimental conditions. The primary distinction is the controller's unit of action. The process-level interventions remove the explicit goal network or adaptive process scheduling while keeping delegated role execution. \condSingleWriter\ removes delegation but also changes the control unit, so it does not isolate process separation.}
\label{tab:experimental-conditions}
\end{table*}
